\section{Glossary}
\label{sec:glossary}

\noindent\emph{Status: a reading aid, not a source. Each entry fixes one sense of a term and points at the section
where that sense is defined and, where there is one, at the section where the quantity is measured. Where an entry
and a section disagree, the section governs: this list is not a place to record a new result.}

Several words here carry a narrower technical sense than their everyday one, ``exact'' being the obvious case, and
a few numbers acquired names that drifted away from the quantities they were computed from. The entries below fix
the senses in the order the terms are met, from Section~\ref{sec:theory} to Section~\ref{sec:worked}, and each
carries the distinction that the corresponding section had to keep straight. An entry that ends with a trap is a
mistake that was made in this report and then corrected, listed so that it is not made twice.

\begin{description}[leftmargin=1.6em,style=nextline,itemsep=4pt]

  \item[Finite integration technique (FIT).]
        A discretisation of Maxwell's equations on a pair of interlocked cell complexes, primal and dual, in which
        the curl and divergence operators are integer incidence matrices and the metric enters only through the
        diagonal Hodge star matrices. What is exact in it is the topology (Proposition~\ref{prop:identities}); what
        is approximated is the constitutive relation and the time step. It is not an exact method, and this report
        does not call it one.

  \item[Primal and dual complex; degrees of freedom.]
        The primal complex is the hexahedral mesh of edges and faces; the dual complex is built on the dual cells
        whose vertices are the primal cell centres. The unknowns are line integrals of $\vec e$ along primal edges
        and of $\vec h$ along dual edges, so that a primal edge carries a voltage and a dual edge a magnetic
        voltage, and one degree of freedom belongs to each. Written bold lower case, with a tilde on a dual
        quantity; the notation table is Table~\ref{tab:notation}.

  \item[Topological operators and the exact identities.]
        The curl matrices $\Cmat$ (primal edges to primal faces) and $\Smat$ (dual edges to dual faces) and the
        gradient $\Gmat$ are sparse, integer, and satisfy $\Smat\Cmat=0$ identically. The discrete identities of
        Proposition~\ref{prop:identities} and \eqref{eq:identities} are consequences of the incidence matrices
        alone: they hold for any material, on any grid, at any time step. This is the part of the scheme that
        cannot be wrong, and it is the reason statements about topology and statements about accuracy are kept
        apart throughout.

  \item[Hodge star matrices.]
        $\Em$, $\Mm$ and $\Ms$: diagonal, positive, with entries assembled from cell measures and material
        constants, as in \eqref{eq:stars} and Section~\ref{sec:hodge}. They are the only spatial approximation in
        the scheme -- diagonal, measure-based, low order -- and every accuracy statement in this report is
        ultimately a statement about them.

  \item[Leapfrog update and the mixed-state invariant.]
        The electric and magnetic unknowns are staggered by half a step (\eqref{eq:leapfrog-h},
        \eqref{eq:leapfrog-e}), which is what makes the step explicit. The quadratic form \eqref{eq:invariant} is
        conserved to rounding (Sections~\ref{sec:invariant} and~\ref{sec:evidence}). Trap: the conserved form is a
        mixed-state quantity, in which the two fields are read half a step apart; it is not the instantaneous
        physical energy, which oscillates at twice the drive frequency even though the invariant does not move.

  \item[CFL bound.]
        The largest time step for which the explicit recursion is stable, \eqref{eq:cfl}, tabulated in
        Table~\ref{tab:cfl} against measured runs. It is set by the smallest cell anywhere in the domain, not by the
        average one. Trap: the bound is a statement about stability and not about accuracy, so stepping well below
        it does not remove the dispersion of Section~\ref{sec:dispersion}.

  \item[Numerical dispersion; grid-dependent phase velocity.]
        The dispersion relation \eqref{eq:dispersion} is exact for the scheme and gives the phase velocity through
        \eqref{eq:phase}, which falls below $c_0$ for every non-zero wave vector (Section~\ref{sec:dispersion}).
        The phase velocity is a property of the grid and the wavelength together, not of the medium, and on this
        grid it is measured per run -- $0.904\,c_0$ for the tracked front of the demonstration -- rather than
        assumed from the permittivity.

  \item[PEC and PMC as constraints on the unknown space.]
        A perfect electric or magnetic conductor is imposed by constraining the tangential unknowns on the
        boundary face, \eqref{eq:pec}; it is not a material and involves no approximation of its own. The cost is
        paid in degrees of freedom, which is why the boundary conditions of this report are described as
        constraints and can be argued about exactly.

  \item[Mur first-order absorbing condition.]
        A one-way condition, first order in the cell size, applied to the six outer planes; implemented in
        \texttt{abc.py} and measured in Section~\ref{sec:res-abc}, \eqref{eq:mur}, at about $-50$\,dB of
        reflection at normal incidence in two dimensions. It is one member of a family of local boundary
        conditions and is not ``the'' absorbing boundary of the report.

  \item[PML; matched lossy layer.]
        A graded artificial conductivity added in a shell around the domain, chosen so that the interface
        reflects little and the wave decays inside the shell; implemented in \texttt{pml.py}, measured in
        Section~\ref{sec:res-pml}, where it absorbs more deeply than the one-way condition does on the same
        geometry. Trap: neither mechanism is \emph{the} boundary of the report. Both are additions to the same
        array, both were measured at normal incidence only, and the sections that measure them
        (Sections~\ref{sec:res-abc} and~\ref{sec:res-pml}) are cited side by side rather than ranked.

  \item[Cavity eigenproblem and the gradient null space.]
        The resonance problem is the generalised eigenproblem \eqref{eq:geneig} on the interior edges, and it is
        singular: every gradient of a scalar potential is a null vector of the curl-curl operator at zero
        frequency. Those modes are physical only as the static limit and are removed by projection before the
        solve (Section~\ref{sec:eigen}). Trap: a small eigenvalue returned with no check against the null space is
        more likely to be a gradient mode than a resonance.

  \item[Port; wave port.]
        A port is the surface through which a known field enters and leaves the domain and on which the
        travelling-wave amplitudes are measured (Section~\ref{sec:sp-definition}). A wave port is one that
        carries the transverse mode profile of the line rather than a lumped terminal pair, which is what allows
        $a$ and $b$ to be defined on it and a power to be assigned to each.

  \item[Mode normalisation.]
        The transverse mode profile $\mathbf{e}_{\text{mode}}$ is scaled so that the mode carries $1$\,W, which
        for the line measured here fixes a voltage scale without the report printing a value for it. The port
        voltage and current are then projections onto that profile, $V=\int\mathbf{E}\cdot
        \mathbf{e}_{\text{mode}}$ and $I=\oint\mathbf{H}\cdot\mathbf{h}_{\text{mode}}$, and the reference
        impedance follows from them (Sections~\ref{sec:sp-definition} and~\ref{sec:sp-power}).

  \item[Travelling waves $a$ and $b$.]
        The incident and outgoing wave amplitudes at a reference plane, normalised so that $|a|^2$ and $|b|^2$
        are powers and their difference is the power crossing the plane, \eqref{eq:sp-power}. They are
        coefficients of a mode expansion on the port surface, not voltages and not field values at a point.

  \item[Reference plane.]
        The transverse surface at which $a$ and $b$ are evaluated and at which the phase of a scattering
        parameter is defined (Section~\ref{sec:sp-refplane}, \eqref{eq:sp-refplane}). Moving it rotates the phase
        and leaves the magnitude alone. Trap: a scattering parameter quoted without its reference plane has a
        phase that does not mean anything, and the factor of two between a plane shift expressed in the line and
        one expressed in its round trip is a mistake this report made once.

  \item[Scattering parameters and reference impedance.]
        $S_{mn}=b_m/a_n$ with every port terminated in a reference impedance $Z_0$. On the line measured here that
        impedance is complex rather than real -- $z_0 = 49.761 + 13.592\,j$\,\si{\ohm}, of magnitude
        \SI{51.584}{\ohm} at \SI{15.278}{\degree}, which is the fit of \eqref{eq:worked-z0} rounded to three
        decimal places -- and that is why every $S$-parameter derived from the record has to be renormalised
        before it is compared with another tool: the reflection coefficient, and so every element of the matrix,
        depends on which reference impedance is assumed
        (Sections~\ref{sec:sp-refz} and~\ref{sec:worked}). Trap: an $S$-parameter whose reference impedance is
        not stated is not a number, because the same physical termination renormalises to different values, and
        the choice is not cosmetic.

  \item[Renormalisation.]
        The closed-form mapping that carries a scattering matrix from one reference impedance to another,
        \eqref{eq:sp-renorm}, \eqref{eq:sp-renorm-2p}, \eqref{eq:sp-renorm-delta} and \eqref{eq:sp-renorm-1p},
        written in terms of the reflection coefficient $\rho$ of the new reference impedance in the old one and
        checked against the series impedance of \eqref{eq:sp-series-z}. It changes magnitudes as well as phases,
        which is what distinguishes it from a reference-plane shift; the two are often confused because both are
        called corrections.

  \item[De-embedding; reference-plane shift.]
        Moving a measured network from the port plane to another plane along the line, by rotating with the
        propagation constant of the line (Sections~\ref{sec:sp-refplane} and~\ref{sec:ms-deembed}). It is exact
        given $\beta$ and it cannot repair an error in magnitude: whatever is wrong with the measurement at the
        port plane is still there, rotated, at the new one.

  \item[Touchstone; \texttt{.s2p}.]
        The plain-text interchange format for scattering parameters: an option line such as
        \texttt{\# HZ S RI R 50}, then one row per frequency, frequency first. The port count is encoded in the
        filename by convention, and this project's reader infers it from the column count instead
        (Section~\ref{sec_touchstone}, \texttt{emfit/ports.py:1849}). Trap: the row order Touchstone v1 mandates
        is input port slowest, so a two-port row reads $S_{11}\,S_{21}\,S_{12}\,S_{22}$; a matrix flattened in the
        natural order transposes the file.

  \item[Passivity.]
        The statement that a network cannot deliver more power than it receives. As a test on what was measured it
        is the circuit condition $|s_{21}|^2+|\gamma|^2\le1$ (\eqref{eq:worked-passivity}), which for the rung 8
        record reads $0.96228$, a loss of $3.77$ per cent. Trap: that percentage is a statement about the circuit
        the measurement defined, and it is not comparable with the flux deficits of the next two entries, which
        are statements about the field at a plane; the two live on different records and must not be subtracted
        from one another.

  \item[Station-profile tilt diagnostic; insertion loss.]
        The end-to-end change in the station Poynting profile on the CW run at \SI{5.8278}{\giga\hertz}, nine
        stations at cells $52$--$100$ of the propagation axis over a \SI{9.6}{\milli\metre} span, $0.987038$,
        which the ledger of Section~\ref{sec:worked} also quotes as $-0.1133$\,dB over that span, the logarithm
        taken in the amplitude convention, $20\log_{10}$. That span is $0.35\,\lambda_g$ at the drive frequency,
        against a standing wave whose ripple period is $\lambda_g/2=\SI{13.688}{\milli\metre}$
        (\eqref{eq:tline-ripple-period}, Sections~\ref{sec:worked-tilt} and~\ref{sec:msres-alpha}). Trap: a
        window shorter than the ripple period cannot supply a decay per unit length, so this number is a
        diagnostic of the profile and not an attenuation, and where earlier sections called it an interim
        insertion loss they were naming a quantity it cannot supply. A reader converting it himself with
        $10\log_{10}$ gets half the figure, which is why the sentence stating the convention sits beside
        the number in Section~\ref{sec:worked}.

  \item[Power deficit and amplitude deficit.]
        One record, two numbers, a factor of two apart. The power ratio $s_9/s_1=0.936635$ gives a deficit of
        $6.3365$ per cent, printed as $6.34$; the amplitude ratio $\sqrt{s_9/s_1}=0.96780$ gives $3.2201$ per
        cent, printed as $3.22$ (Sections~\ref{sec:worked-ledger} and~\ref{sec:worked-tilt}). Trap: a deficit
        quoted as a bare percentage is ambiguous until it says which ratio it came from, and a deficit is
        compared with a tilt only when both are referenced the same way.

  \item[$\eps_{\text{eff}}$ and $\lambda_g$.]
        The effective permittivity of the quasi-TEM line and the guided wavelength that follows from it:
        $\eps_{\text{eff}}=3.5322$ gives $\lambda_g=\SI{27.375}{\milli\metre}$ against \SI{51.449}{\milli\metre}
        in free space at the same frequency (Sections~\ref{sec:tline-epseff} and~\ref{sec:msres-import}). Trap:
        the closed-form estimate differs by $7.8$ per cent and that bias belongs to the reference box, not to the
        solver, so agreement with closed form is not by itself the test of either.

  \item[Accumulation window and warm-up.]
        The steady-state interval over which the frequency-domain sums are taken: of $2400$ steps, $600$ were
        discarded and $1800$ accumulated, a window of $6.8637\times10^{-10}$\,s (Section~\ref{sec:worked}). Trap:
        quantities read off a plane are raw sums over that window while station quantities are means of the
        Poynting flux, so the two are not comparable until each is divided by the window length, and a factor of
        the window length appearing between two columns is a bookkeeping error and not physics.

\end{description}

The list stops here because the terms beyond it are used once and defined where they are used. It is deliberately
short of formulas: where a formula is written above it is a pointer to the numbered equation that carries the
argument, and the equation, not the entry, is the statement to be checked.
